\chapter{Discussion}
\label{ch:discussion}

This chapter interprets the results presented in Chapter~4 and relates them to
the project objective of developing an efficient and scalable computer-vision
workflow for robot interaction with elevator infrastructure on embedded
hardware.

\section{Interpretation of the main findings}

The main outcome of this project is that a practical embedded perception
workflow can be achieved by combining instance-segmentation-based button
detection, crop-based \gls{ocr}, and a modular \gls{ros2} runtime pipeline. The
results show that the strongest overall system did not come from selecting the
newest model architecture or from maximizing validation performance in isolation.
Instead, the most effective workflow emerged from a sequence of engineering
decisions that balanced perception quality, deployment constraints, and future
scalability.

For \gls{yolo}, the baseline \gls{yolo}v8-Seg model remained the most reliable
choice for deployment. Hyperparameter tuning improved validation performance, but
those gains did not translate into a better final held-out test model. Likewise,
YOLO26-small did not outperform the baseline under the same evaluation
conditions. This is an important result because it shows that, for this task, a
carefully chosen and well-understood baseline can be more valuable than adopting
a newer architecture without a clear deployment benefit. The most practically
useful \gls{yolo} improvement came from environment-specific fine-tuning with
\gls{adc} data from HAN University, Building 31, which substantially improved
performance in the target environment while preserving the original-domain
performance.

For \gls{ocr}, the project produced a clearer tuning outcome. The selected stage-2
configuration improved both test accuracy and normalized edit distance relative
to the initial recognizer. In contrast to the \gls{yolo} tuning results, the
\gls{ocr} validation improvements did translate into better held-out test-set
performance. This suggests that the \gls{ocr} stage was more sensitive to
optimization settings in a way that remained meaningful at test time.

At system level, the deployed \gls{ros2} pipeline demonstrates that model
quality alone is not sufficient for real-time embedded performance. The final
pipeline throughput depended strongly on how much non-inference work remained in
the runtime path. The shift from full-image mask export to compact ROI-local mask
handling, together with the split-node \gls{yolo}/\gls{ocr} architecture and the
bounded \gls{ocr} workload, produced the most important runtime gains. The best
measured configuration achieved about 13.1~\gls{fps}, while the collector-enabled
configuration remained practical at about 10~\gls{fps}. This shows that the
embedded performance target was not reached by model export alone, but by a
combination of model export, message-design decisions, and workload placement.

\section{Relation to the research questions}

The results provide a coherent answer to the main research question. An
efficient and scalable computer-vision workflow for hospital interaction-point
perception can be designed by using instance segmentation for button localization,
crop-based \gls{ocr} for semantic interpretation, \gls{tensorrt} export for
embedded execution, and an \gls{ros2}-based split-node architecture for modular
runtime integration.

The first sub-question concerned the relevant human-interaction points in
hospital environments. The project focused on elevator buttons and related panel
interfaces because they are a concrete and operationally critical interaction
class for multi-floor service-robot mobility. This focus proved appropriate: the
task is sufficiently challenging to expose problems of detection, segmentation,
recognition, and embedded deployment, while still being narrow enough to support
a complete workflow.

The second sub-question asked which methods and architectures were suitable.
The literature review suggested that instance segmentation is better aligned with
precise robot interaction than coarse object detection alone, and the results
support this choice. The segmentation masks were directly useful for masked depth
estimation and 3D target generation. The final architecture therefore combined a
segmentation-based detector with a separate recognizer rather than relying on a
single end-to-end model.

The third sub-question concerned dataset design. The results show that dataset
quality and dataset relevance were both important. The original curated dataset
was sufficient to train a strong baseline, but the environment-specific
fine-tuning experiment demonstrated that even a relatively small amount of
deployment-collected data can substantially improve performance in the target
domain. This supports the decision to include an \gls{adc} in the full workflow.

The fourth sub-question asked which optimization techniques support embedded
inference. Here the answer is twofold. First, \gls{tensorrt} export was
beneficial for both \gls{yolo} and \gls{ocr}, with especially strong gains for
\gls{ocr}. Second, export alone was not enough: runtime-side optimizations such
as reducing annotation overhead, separating the inference nodes, and limiting the
number of processed \gls{ocr} regions were equally important.

The fifth sub-question concerned scalability and efficient future improvement.
The final system addresses this directly through the \gls{adc}, which turns
runtime use into a source of curated retraining data. The environment-specific
\gls{yolo} improvement demonstrates that this design is not only theoretically
scalable, but also practically useful.

\section{Limitations}

Several limitations should be considered when interpreting the results. First,
the project focused on elevator-interface perception rather than the full set of
hospital interaction elements. While this was a justified and useful scope
reduction, the final conclusions should not automatically be generalized to all
door panels, access readers, or human-machine interfaces.

Second, some benchmark results remain sensitive to the evaluation setup. For
example, the single-image framework-to-\gls{tensorrt} comparison is useful for
illustrating acceleration trends, but it is not by itself a complete deployment
benchmark. Likewise, some distance-related \gls{yolo} improvements were most
clearly observed in field behavior and specific replay analyses rather than in a
single universal summary metric. For this reason, the confidence and coverage
improvements are stronger conclusions than any claim of a general doubling of
recognition range.

Third, the project was executed on a single embedded platform and camera setup.
The \gls{jetson} Orin Nano and Stereolabs ZED X are representative for the
target use case, but performance and tuning decisions may differ on other edge
devices or other sensing configurations.

Fourth, the report currently emphasizes perception and runtime integration rather
than downstream manipulation or actuation. The system produces 3D button targets
and button labels, but the full closed-loop robot interaction behavior lies
beyond the present project scope.

\section{Implications for future robot deployment}

Despite these limitations, the project has clear practical implications. The
results suggest that a service robot does not require a monolithic perception
stack to achieve useful embedded performance. A modular pipeline with explicit
message interfaces can be easier to tune, benchmark, and extend. This is
particularly important for educational and research robots such as CareBot, where
iterative development and subsystem replacement are expected.

The project also suggests that deployment data should be treated as a core part
of the system design rather than as an afterthought. The inclusion of the
\gls{adc} makes the perception workflow self-improving in a structured way: the
same runtime system used in deployment can collect higher-value examples for the
next training cycle. This closes an important loop between experimentation,
deployment, and future model improvement.

Overall, the project shows that the combination of instance segmentation,
task-specific \gls{ocr}, \gls{tensorrt} acceleration, and structured deployment
data collection provides a credible foundation for later integration into a more
complete autonomous hospital robot.
